% Replace the current "Within-participant analysis." and
% "Between-participant analysis." paragraphs in §4.5 with this single paragraph.

\textbf{Tempo condition analysis.}
Cadence (steps per minute) is derived from heel-strike peaks in the pressure
signal and evaluated on a rolling 1\,Hz grid.
Each participant's natural walking pace is estimated individually from buffer
windows at the boundary of their sessions, scaled to their own baseline rather
than a fixed population mean.
The trial is divided into three phases aligned to the audio ramp structure ---
ramp-in (0--20\%), plateau (20--80\%), and ramp-out (80--100\%) --- and cadence
deviations in each phase are expressed as z-scores relative to that participant's
own cadence variability during the weight condition, making individual responses
comparable regardless of stride consistency.
Sessions are assigned one of 11 behavioural profiles based on the shape of the
cadence trajectory, and a continuous influence score ($s \in [0,1]$) summarises
the magnitude of the largest detected response across any phase.
The same pipeline is applied to weight-condition data as a false-positive
benchmark.
Group-level cadence shift is tested with the Wilcoxon signed-rank test within
each direction subgroup; Spearman's rank correlation between influence score and
Agency Aggregate provides the primary test of Hypothesis~1.
Full algorithmic detail, feature definitions, classification criteria, and
edge-case handling are provided in Appendix~\ref{app:tempo-pipeline}.
